\documentclass[12pt]{article}
\usepackage[utf8]{inputenc}
\usepackage[russian]{babel}
\usepackage{amsmath}
\usepackage{amssymb}
\usepackage{geometry}
\geometry{a4paper, margin=2cm}

\begin{document}

\section*{Задача 3.1.5}

\textbf{Условие:} Определить, является ли последовательность \(a_1, \dots, a_n\) отсортированной в порядке возрастания элементов, то есть \(a_1 \le \dots \le a_n\). Вывести 1, если является, или 0, если не является (0.5 балла):\\
in: n, a1, ..., an \\
out: r in \{0,1\}

\section*{Решение}

\begin{verbatim}
Read 1  ; n
Load =0
Sub 1
JZero SUCCESS   ; n = 0
JGTZ SUCCESS    ; n < 0

Read 2  ; a_i

Load 1
Sub =1
Store 1 ; n := n-1

WHILE: Load 1
    JZero SUCCESS
    Read 3  ; r3 := a_{i+1}
    Load 2
    Sub 3   ; r0 := a_i - a_{i+1}
    JGTZ FAIL
    Load 1
    Sub =1
    Store 1 ; n := n-1
    Load 3
    Store 2 ; r2 := a_{i+1}
    Jump WHILE
SUCCESS: Write =1
    Halt 0
FAIL: Write =0
    Halt 0
\end{verbatim}

\section*{Сложность}

Худший случай - уже отсортированная по возрастанию последовательность длины $n > 0$, так как тогда выполняется максимальное количество сравнений.

\begin{enumerate}
    \item Инициализация выполняется ровно за 9 шагов
    \item Цикл WHILE для проверки пар элементов выполняется n-1 раз
    \item На каждой успешной итерации выполняется 12 шагов
    \item На n-й итерации выполняется проверка условия - 2 шага
    \item SUCCESS занимает 2 шага
\end{enumerate}

Значит сложность в худшем случае:
\[
    T_{\max}(n) = 9 + 12(n - 1) + 2 + 2 = 12n + 1
\]

\end{document}